\chapter{Introduction} \label{chapter:introduction}

\label{para:interface_importance} % ---------- %

Material interfaces underpin the performance and reliability of advanced solid-state systems.

In semiconductor devices, they govern charge transport, carrier confinement, defect formation, and phase stability; in energy conversion and storage technologies, they regulate ion exchange kinetics, chemical compatibility, and thermal management.

In emerging heterostructures, interfacial interactions give rise to novel phenomena, ranging from encapsulated metastable phases to Schottky barrier formation that are absent in the bulk constituents.

As device dimensions shrink and surface-to-volume ratios increase, these interfacial effects become ever more pronounced, highlighting the critical need for accurate prediction and control of interface structure and energetics.

\label{para:interface_foundation} % ---------- %

Material interfaces delineate the boundary where two distinct solids meet, whether as heterostructure junctions, grain
boundaries, or thin films, and exhibit unique atomic, electronic, and thermal characteristics that set them apart from
the bulk phases.

Local features such as strain fields, atomic reconstruction, and elemental intermixing create modified bonding
environments and electronic states, profoundly influencing the mechanical, transport, and optical behaviour of the
composite.

In today's micro- and nanoscale devices, these interfacial phenomena dominate performance as components shrink, making
rigorous characterisation and predictive modelling of interfaces indispensable for materials design.

\label{para:material_interface_definition} % ---------- %

A material interface is the region where two distinct solids meet and form a junction, commonly arising in
heterostructures, grain boundaries, or thin film systems.

The local structure at the interface can include strain fields, atomic reconstruction, and intermixing of constituent
elements, all of which can significantly influence the behaviour of the composite system.

In this research, material interfaces are studied predominantly through \textit{ab initio} methods, particularly Density
Functional Theory (DFT), with a focus on predicting and analysing the structural and energetic properties of 2D|2D and
2D|3D junctions.

Interfaces are modelled as systems with finite thickness that bridge two crystalline lattices.

Depending on their geometric orientation, interfaces are categorised as either \emph{stacked}, where a 2D layer rests
atop a substrate, or \emph{lateral}, where edge-on growth leads to in-plane bonding between dissimilar materials.

Band offsets at interfaces, often deviating from predictions by simple models like Anderson's rule, are governed by
local charge redistribution, dipole formation, and interface-specific states.

The emergence of new electronic phases at interfaces, such as Ohmic contact formation or type-II band alignments,
necessitates accurate modelling beyond bulk considerations.

Several methodologies have been developed to address the challenges of constructing realistic interface models.

The ARTEMIS tool enables high-throughput generation of interface structures through lattice matching and surface
termination selection, allowing systematic exploration of the complex interfacial phase space.

Similarly, the RAFFLE method uses statistical learning and void-filling algorithms to construct interface phases that
are energetically favourable and representative of plausible morphologies.

Interfaces also significantly influence phonon transport.

They can reduce thermal conductivity compared to bulk values; an effect exploitable in thermoelectric applications.

Interfacial scattering, acoustic impedance mismatch, and atomic-scale disorder all contribute to this behaviour,
reinforcing the interface's role as a functional component rather than a passive boundary.

Thus, in the context of this research, a material interface is not merely a geometric boundary but an emergent,
functionally distinct region.

It governs critical physical properties, demands precise structural modelling, and serves as a fertile ground for
material innovation and device engineering.

\label{para:interface_classification} % ---------- %

Material interfaces can be broadly categorised by their crystallographic orientation, dimensionality, and degree of
lattice coherence.

These classifications are instrumental in understanding interfacial phenomena and in designing simulation protocols that
are physically realistic and computationally tractable.

\label{para:interface_coherence} % ---------- %

\paragraph{By coherence: coherent, semi-coherent, incoherent.} Interfaces can also be distinguished by the degree of
lattice match across the boundary:

\begin{itemize}

    \item \textbf{Coherent interfaces} exhibit perfect registry between atomic planes, typically occurring when lattice
    mismatch is negligible.

    These interfaces preserve the translational symmetry across the junction and are often found in homoepitaxial or
    pseudomorphic systems.

    \item \textbf{Semi-coherent interfaces} accommodate lattice mismatch through periodic misfit dislocations.

    These dislocations relieve interfacial strain while maintaining partial registry, allowing larger domains of
    coherency interspersed with localised defects.

    \item \textbf{Incoherent interfaces} display no periodicity match across the interface, leading to structurally
    disordered junctions.

    These are often found in polycrystalline or amorphous/crystalline composites and can act as strong phonon or
    electron scattering centres.

\end{itemize}

\label{para:chemical_contrast_hetero_homo} % ---------- %

By chemical contrast: heterojunctions vs homojunctions.

Chemical composition also plays a defining role.

\emph{Homojunctions} are interfaces within a single material where the crystal lattice remains continuous, but local
properties differ, such as variations in dopant concentration or applied strain, resulting in regions with distinct
electronic or mechanical behaviour, while \emph{Heterojunctions} involve distinct materials with differing
electronegativity, band gaps, and ionic/covalent character.

Heterojunctions are central to semiconductor devices, solar cells, and catalysis, often giving rise to emergent
interfacial dipoles, states, or reconstructed bonding motifs.

\label{para:interface_classifications_continued} % ---------- %

These classifications are not mutually exclusive; real interfaces may be simultaneously lateral and semi-coherent, or
stacked but chemically graded.

Understanding these distinctions is essential when generating or interpreting simulation results for both idealised and
realistic interfacial systems.

\label{para:interface_effects_2D_oxides} % ---------- %

Material interfaces influence a wide array of physical and electronic properties, often determining the functional
behaviour of composite systems and devices.

These effects are particularly pronounced in systems involving 2D materials, complex oxides, or strongly correlated
electron systems.

\label{para:band_alignment_interfaces} % ---------- %

Perhaps the most critical electronic property influenced by interfaces is the band alignment between materials.
Misalignment of conduction and valence band edges creates energy barriers or wells that determine whether an interface
behaves as a Schottky barrier, Ohmic contact, or heterojunction.

These alignments are crucial for device operation in field-effect transistors, diodes, and photovoltaics.
Interfacial dipoles, charge transfer, and electronic reconstruction can shift the effective band edges, often
invalidating simple rules like Anderson's model.

\label{para:interface_states} % ---------- %

Interfaces often host localised states arising from atomic mismatch, bonding reconstruction, or impurity segregation.
These can act as traps or recombination centres, significantly affecting carrier lifetimes and mobility.

\label{para:dielectric_interface_permittivity} % ---------- %

The dielectric properties of materials can be strongly altered at interfaces.

In systems like \ce{CaCu3Ti4O12}, colossal permittivity arises not from the bulk crystal but from insulating grain
boundaries and a semiconducting interior, exemplifying the Maxwell-Wagner effect.

Interface-induced polarisation, charge accumulation, and ionic displacements all contribute to enhanced dielectric
responses.

\label{para:thermal_transport_interfaces} % ---------- %

Thermal transport across interfaces is typically suppressed due to phonon scattering, mass disorder, and acoustic
impedance mismatch.

This results in thermal boundary resistance (or Kapitza resistance), a key factor in thermoelectric design and heat
management in microelectronics.

In 2D heterostructures, the interface often acts as a dominant phonon scattering site, reducing the effective thermal
conductivity below that of either constituent layer.

\label{para:interface_stabilise_metastable_phases} % ---------- %

Interfaces can stabilise new or metastable phases not accessible in bulk.

For instance, the interface between \ce{BaTiO3} and \ce{SiO2} has been shown to promote the spontaneous formation of the
fresnoite phase \ce{Ba2TiSi2O8}, altering both electronic and mechanical response.

A more nuanced example is the graphene|\ce{MgO} interface: while monolayer MgO confined between graphene layers often
adopts a rocksalt-like structure, this is not simply a case of one phase being energetically favoured.

Instead, the local environment, including charge transfer and encapsulation, stabilises the rocksalt phase relative to
the expected hexagonal phase in unsupported monolayers.

However, experimental and theoretical studies suggest that a mixed or hybrid phase may also emerge with phase stability
dependent on thickness and interface strain.

\label{para:interface_mechanical_properties} % ---------- %

Mechanical properties such as bulk modulus, shear strength, and adhesion energy are modified at interfaces due to
altered bonding environments and lattice mismatch.

These variations are important in the design of composite materials, flexible electronics, and nanoindentation studies.

\label{para:interface_optical_effects} % ---------- %

Interfaces affect the optical absorption spectrum by introducing new electronic states and modifying the joint density
of states.

This can lead to bandgap narrowing or broadening, shifts in refractive index, and enhancement of excitonic effects;
features exploited in photodetectors and photovoltaic devices.

\label{para:interface_summary} % ---------- %

Material interfaces act as multifunctional domains that modulate a broad spectrum of material
properties.

Their control and prediction are vital for the design of next-generation electronic, optoelectronic, and thermal
devices.

\label{para:interface_functional_regions} % ---------- %

Interfaces are not merely passive boundaries between materials; they have an active impact on their surrounding
environment, functional regions that frequently control and even dominate the physical properties of the system.

As materials and devices are engineered at ever smaller scales, the role of interfaces becomes more pronounced.

In conventional bulk systems, interfaces may be limited to grain boundaries or inter-phase regions.

In nanomaterials, layered heterostructures, and device-scale systems, interfacial regions can constitute a
non-negligible fraction of the total volume or surface area.

Their influence is therefore not ancillary but foundational to system performance.

\label{para:interface_stability_challenge} % ---------- %

Yet despite their importance, predicting which interfaces are energetically favourable remains a significant challenge.

Interface stability is not solely determined by bulk lattice matching, but also by local bonding, chemical
inhomogeneity, and atomic-scale reconstruction.

The structural space is vast: each pair of surfaces can be stacked with arbitrary lateral shifts,terminations, and
alignments, resulting in a combinatorial explosion of configurations.

Strain relaxation, interlayer interactions, and charge transfer further complicate the energy landscape.

Accurately assessing favourability across this space demands a balance between resolution and scale that conventional
methods struggle to achieve.

\label{para:dft_scalability_bottleneck} % ---------- %

First-principles techniques such as Density Functional Theory (DFT) remain the standard for computing accurate
interfacial energies and relaxations.

However, their high computational cost, scaling cubically with system size, renders exhaustive sampling intractable for
realistic sizes structures, such as a 22 \times Si to 21 \times Ge latticed matched supercell, which often comprise
hundreds of atoms.

This bottleneck precludes broad screening efforts, especially when exploring multiple chemistries,orientations, and
stacking arrangements.

As such, DFT might be better suited to final validation and a confirmation of datasets evaluated in a less
computationally expensive manner.

\label{para:mlp_screening} % ---------- %

To address this, the present project explores the use of machine-learned interatomic potentials (MLPs) for scalable
interface screening.

MLPs are data-driven models trained to reproduce DFT-level energies and forces, but at orders of magnitude lower cost.

In particular, this work focuses on the MACE framework; an equivariant message-passing architecture that preserves
rotational symmetries and captures local atomic interactions with high fidelity.

By using MACE to rapidly evaluate large numbers of candidate interfaces, it becomes possible to rank structures by
predicted stability and prioritise low energetic configurations for subsequent DFT validation.

\label{para:hierarchical_modelling_workflow} % ---------- %

This approach supports a hierarchical modelling workflow in which DFT and MLPs are used in tandem: MLPs for breadth, and
DFT for depth.

It enables rigorous testing of whether MLPs can replicate DFT-derived favourability rankings across a diverse set of
semiconductor systems, including phase boundaries (e.g. Si|Si) and heterostructures (e.g. Si|Ge, Ge|C).

For the context of this study, I ignore the semicontuctor requirement in its definition.

Moreover, it lays the foundation for further improvements through techniques such as delta-learning, structure-based
prediction models, and surface-to-interface generalisation.

\label{para:report_structure} % ---------- %

The remainder of this report is structured as follows:

\begin{itemize}

    \item \textbf{Chapter~\ref{chapter:background}} introduces the theoretical foundations and relevant literature,
    covering interface energetics, DFT formalism, and the development of MLPs.

    \item \textbf{Chapter~\ref{chapter:computational_investigation}} presents a benchmarking study comparing MLP and DFT
    predictions of interface favourability, including statistical correlation analysis and system-specific case studies.

    \item \textbf{Chapter~\ref{chapter:next_steps}} outlines planned extensions to the current methodology, including
    the development of delta-corrected potentials and predictive models linking surface properties to interface
    behaviour.

    \item \textbf{Chapter~\ref{chapter:conclusions}} concludes with a discussion of findings, limitations, and broader
    implications for autonomous materials modelling and the design of functional heterostructures.

\end{itemize}

\label{para:framework_objective} % ---------- %

Through this integrated modelling framework, the project seeks to reduce the cost and uncertainty of interface
prediction while preserving the fidelity needed for scientific insight and technological relevance.